\chapter{Stochastic Differential Equation}

Consider a stochastic process satisfying:
\begin{equation*}
    dX_t = b(t,X_t,W_t)dt + \sigma(t,X_t,W_t)dW_t \tag{$*$}.
\end{equation*}
for two problems:
\begin{enumerate}[label=(\arabic{*})]
    \item existence and uniqueness and properties of solutions,
    \item how to solve for particular cases.
\end{enumerate}

\section{Examples}

\begin{exam}[Geometric Brownian Motion]
    Solving
    \begin{equation*}
        dX_t = \alpha X_t dt + \sigma X_t dW_t.
    \end{equation*}
    where the initial value $X_0$ is given, and $\alpha,\sigma$ are constant.

    \noindent \emph{Solution:} Diving $X_t$ on the both side and integrating,
    \begin{equation*}
        \int_0^t \frac{dX_u}{X_u} = \alpha t + \sigma W_t.
    \end{equation*}
    Let $f(x) = \log x$. Then by It\^o formula,
    \begin{align*}
        d \log X_t &= \frac{1}{X_t}dX_t - \frac{1}{2} \frac{1}{X_t^2} (dX_t)^2 \\
        &= \frac{1}{X_t}dX_t - \frac{1}{2} \frac{1}{X_t^2} \sigma^2X_t^2 dt\\
        &= \frac{1}{X_t}dX_t - \frac{1}{2} \sigma^2 dt.
    \end{align*} 
    So
    \begin{equation*}
        \log X_t - \log X_0 = \int_0^t \frac{dX_u}{X_u} - \frac{1}{2} \sigma^2 t.
    \end{equation*}
    It follows that
    \begin{equation*}
        X_t = X_0\exp\bc{\sigma W_t + \bc{\alpha - \frac{1}{2}\sigma^2}t}.
    \end{equation*}
    \begin{rmk}
        \begin{enumerate}[label=(\arabic{*})]
            \item If $(W_t)_{t \geq 0}$ is independent of $X_0$, then
            \begin{align*}
                \E[X_t] &= \E[X_0]\E\bj{\exp\bc{\sigma W_t + \bc{\alpha - \frac{1}{2}\sigma^2}t}} \\
                &= e^{\alpha t}\E[X_0]\E\bj{\exp\bc{\sigma W_t - \frac{1}{2}\sigma^2t}} \\
                &=e^{\alpha t}\E[X_0],
            \end{align*}
            because $\exp\bc{\sigma W_t - \frac{1}{2}\sigma^2t}$ is a martingale.

            \item If $\alpha > \frac{1}{2}\sigma^2$, then $X_t \sto \infty$ as $t \sto \infty$. If $\alpha < \frac{1}{2}\sigma^2$, $X_t \sto 0$ as $t \sto \infty$. If $\alpha = \frac{1}{2}\sigma^2$, $X_t$ will fluctuate between arbitrary large and arbitrary small values.
        \end{enumerate}
    \end{rmk}
    \begin{defn}
        A stochastic process $(X_t)_{t \geq 0}$ of the form
        \begin{equation*}
            X_t = X_0\exp(\sigma W_t + \mu t)
        \end{equation*}
        is called the geometric Brownian motion.
    \end{defn}
\end{exam}

\begin{exam}[Hull-White Interest Rate Model]
    Consider SDE
    \begin{equation*}
        dR_t = (a_t - b_tR_t) + \sigma_t dW_t,\quad R_0 = r,
    \end{equation*}
    where $a_t,b_t,\sigma_t$ are deterministic.

    \noindent \emph{Solution:} 
    \begin{equation*}
        dR_t + b_tdR_t = a_tdt + \sigma_t dW_t.
    \end{equation*}
    Multiplying $e^{\int_0^t b_u du}$,
    \begin{equation*}
        e^{\int_0^t b_u du}dR_t + e^{\int_0^t b_u du}b_tdR_t = e^{\int_0^t b_u du}a_tdt + \sigma_te^{\int_0^t b_u du} dW_t.
    \end{equation*}
    Because $e^{\int_0^t b_u du}$ is of bounded variation on any interval, by It\^o formula
    \begin{equation*}
        d \bc{e^{\int_0^t b_u du}}dR_t = e^{\int_0^t b_u du}a_tdt + \sigma_te^{\int_0^t b_u du} dW_t.
    \end{equation*}
    So
    \begin{equation*}
        R_t = re^{-\int_0^t b_u du} + r \int_0^t e^{\int_t^s b_u du}a_sds + \int_0^t\sigma_se^{\int_t^s b_u du} dW_s.
    \end{equation*}
\end{exam}

\begin{exam}
    Consider SDE
    \begin{equation*}
        dX_t = rX_t(K- X_t)dt + \beta X_t dW_t,\quad X_0 = x > 0.
    \end{equation*}

    \noindent \emph{Solution:}
    \begin{equation*}
        \frac{1}{X_t}dX_t + rX_tdt =rKdt +  \beta dW_t.
    \end{equation*}
    Therefore,
    \begin{equation*}
        \int_0^t\frac{1}{X_t}dX_t + \int_0^t rX_tdt = rKt + \beta W_t.
    \end{equation*}
    For the left hand side, first
    \begin{align*}
        \int_0^t\frac{1}{X_t}dX_t &= \log X_t -\log X_0 +\frac{1}{2} \int_0^t \frac{1}{X_s^2} (dX_t)^2 \\
        &= \log X_t -\log X_0 +\frac{1}{2} \int_0^t \frac{1}{X_s^2} \beta^2 X_s^2 ds \\
        &= \log \frac{X_t}{x} + \frac{1}{2}\beta^2 t.
    \end{align*}
    Therefore,
    \begin{equation*}
        X_t\exp\bc{r\int_0^tX_sds} = x\exp\bc{\beta W_t + \bc{rK - \frac{1}{2}\beta^2}t}.
    \end{equation*}
    Integrating w.s.t. $t$,
    \begin{align*}
        \int_0^t \exp\bc{r\int_0^tX_sds} d\bc{\int_0^tX_sds} &= \frac{1}{r} \bc{\exp\bc{r\int_0^tX_sds} - 1}\\
        &= x\int_0^t\exp\bc{\beta W_s + \bc{rK - \frac{1}{2}\beta^2}s}ds.
    \end{align*}
    It follows that
    \begin{equation*}
        \exp\bc{r\int_0^tX_sds} = 1+ rx\int_0^t\exp\bc{\beta W_s + \bc{rK - \frac{1}{2}\beta^2}s}ds,
    \end{equation*}
    and so
    \begin{equation*}
        \int_0^tX_sds  = \frac{1}{r}\log \bc{1+ rx\int_0^t\exp\bc{\beta W_s + \bc{rK - \frac{1}{2}\beta^2}s}ds}.
    \end{equation*}
    Differentiating w.s.t. $t$,
    \begin{equation*}
        X_t = \frac{x\exp\bc{\beta W_t + \bc{rK - \frac{1}{2}\beta^2}t}}{1+ rx\int_0^t\exp\bc{\beta W_s + \bc{rK - \frac{1}{2}\beta^2}s}ds}.
    \end{equation*}
\end{exam}

\begin{exam}
    Consider SDE
    \begin{equation*}
        dX_t = \alpha_t dt + b_t X_tdW_t.
    \end{equation*}

    \noindent \emph{Solution:} Trying to find an integrator $\rho_t$ for
    \begin{align*}
        \rho_tdX_t - b_t\rho_t X_tdW_t = \alpha_t\rho_t dt.
    \end{align*}
    By It\^o formula, 
    \begin{equation*}
       d(\rho_t X_t) = \rho_t dX_t + X_t d\rho_t+ (dX_t)(d\rho_t) 
    \end{equation*}
    First, try to find $\rho_t$ such that
    \begin{equation*}
        X_t d\rho_t = - b_t\rho_t X_tdW_t~\Rightarrow~ \frac{d\rho_t}{\rho_t} = -b_tdW_t.
    \end{equation*}
    Then by above
    \begin{equation*}
        \log \frac{\rho_t}{\rho_0} + \frac{1}{2}\int_0^t \frac{1}{\rho_t^2}(d\rho_t)^2 = \log \frac{\rho_t}{\rho_0} + \frac{1}{2}\int_0^tb_u^2du = - \int_0^t b_u dW_u,
    \end{equation*}
    which implies that
    \begin{equation*}
        \rho_t = \exp\bc{- \int_0^t b_u dW_u -  \frac{1}{2}\int_0^tb_u^2du}
    \end{equation*}
    by setting $\rho_0 = 1$. Then
    \begin{align*}
        d(\rho_t X_t) &= \rho_t dX_t + X_t d\rho_t+ (dX_t)(d\rho_t) \\
        &= \rho_t dX_t - - b_t\rho_t X_tdW_t - b_t^2\rho_tX_t dt \\
        &= \alpha_t\rho_t dt - b_t^2\rho_tX_t dt.
    \end{align*}
    So
    \begin{equation*}
        d(\rho_t X_t) +  b_t^2\rho_tX_t dt = \alpha_t\rho_t dt.
    \end{equation*}
    Multiplying $e^{\int_0^t b_u^2 du}$, then we have
    \begin{equation*}
        d\bc{e^{\int_0^t b_u^2 du} \rho_t X_t} = e^{\int_0^t b_u^2 du}\alpha_t\rho_t dt.
    \end{equation*}
    It follows that
    \begin{equation*}
        \rho_t X_t - X_0 = \int_0^t e^{\int_t^s b_u^2 du}\alpha_s\rho_s ds,
    \end{equation*}
    i.e.,
    \begin{equation*}
        X_t = X_0 \exp \left(\int_0^t b_s d W_s-\frac{1}{2} \int_0^t b_s^2 d s\right)+\int_0^t \alpha_s \exp \left(\int_s^t b_u d W_u-\frac{1}{2} \int_s^t b_u^2 d u\right) d s.
    \end{equation*}
\end{exam}

\begin{exam}
    Consider SDE
    \begin{equation*}
        LX_t^{\prime\prime} + RX_t^\prime + \frac{1}{2}X_t = G_t + \alpha \tilde{W}_t,
    \end{equation*}
    where $\tilde{W}_t$ is white noise, i.e., $\tilde{W}_t dt = dW_t$ (in distribution meaning).

    \noindent \emph{Solution:} Introduce
    \begin{equation*}
         \bm{X}_t = \bc{
            \begin{array}{c}
                X_t \\
                X_t^\prime
            \end{array}
        }
    \end{equation*}
    and it follows that
    \begin{equation*}
        d\bm{X}_t = A\bm{X}_t dt + \bm{H}_t dt + \bm{K} dW_t,
    \end{equation*}
    where
    \begin{equation*}
        A=\left(\begin{array}{cc}
            0 & 1 \\
            -\frac{1}{C L} & -\frac{R}{L}
        \end{array}\right),~\bm{H}_t = \bc{
            \begin{array}{c}
                0 \\
                \frac{G_t}{L}
            \end{array}
        },\quad \bm{K} = \bc{
            \begin{array}{c}
                0 \\
                \frac{\alpha}{L}
            \end{array}
        }.
    \end{equation*}
    So
    \begin{equation*}
        \bm{X}_t=\exp (A t) \bm{X}_0+\exp (A t) \int_0^t \exp (-A s)\left(\bm{H}_s d s+\bm{K} d W_s\right).
    \end{equation*}
\end{exam}

\section{Existence and Uniqueness of Solution}

\begin{thm}[Existence and Uniqueness]
    Fix $T > 0$ ($T$ could be $\infty$). Let $b \colon [0,T] \times \R^n \sto \R^n$ and $\sigma \colon [0,T] \times \R^n \sto \R^{n \times m}$ be measurable such that
    \begin{equation}\label{eq:solcond1}
        \norm{b(t,x)} + \norm{\sigma(t,x)}_{\op{F}} \leq K(1 + \norm{x}), \quad \forall~x \in \R^n,~t \in [0,T]
    \end{equation}
    and
    \begin{equation}\label{eq:solcond2}
        \norm{b(t,x)-b(t,y)} + \norm{\sigma(t,x)-\sigma(t,y)}_{\op{F}} \leq K\norm{x-y},\quad \forall~x,y \in \R^n,~t \in [0,T]
    \end{equation}
    for some constant $K$. Let $\bm{Z}$ be a random variable independent of $\mathcal{F}^W_\infty$ and $\bm{Z} \in L^2$. Then
    \begin{equation*}
        d\bm{X}_t = b(t,\bm{X}_t)dt+\sigma(t,\bm{X}_t)d\bm{W}_t,\quad 0\leq t \leq T,~\bm{X}_0 = \bm{Z},
    \end{equation*}
    has a unique (strong) solution $\bm{X}=(\bm{X}_t)_{t \in [0,T]}$ that is continuous w.s.t. $t$ with properties
    \begin{enumerate}[label=(\roman*)]
        \item $\bm{X}_t$ is measurable w.s.t. $\mathcal{F}^W_t \vee \sigma(\bm{Z})$ for all $t$,
        \item $\E\bj{\sup_{0 \leq s \leq T} \norm{\bm{X}_s}^2} < \infty$.
    \end{enumerate}
\end{thm}
\begin{rmk}
    Note that condition (\ref{eq:solcond2}) implies condition (\ref{eq:solcond1}). For the uniqueness, it means that indistinguishable.
\end{rmk}
\begin{proof}
    First, let $\mathcal{E}_T$ be the set of all continuous and $\mathbb{F}^W$-adapted process $\bm{X}$ such that 
    \begin{equation*}
        \E\bj{\sup_{0 \leq s \leq T} \norm{\bm{X}_s}^2} < \infty.
    \end{equation*}
    Then define a norm on $\mathcal{E}_T$ as
    \begin{equation*}
        \norm{\bm{X}}_E^2 \defeq \E\bj{\sup_{0 \leq s \leq T} \norm{\bm{X}_s}^2}.
    \end{equation*}
    It is not hard to see $(\mathcal{E}_T,\norm{\cdot}_E)$ is a Banach space. Note $\bm{X} = \bm{Y}$ in $\mathcal{E}_T$ means that they are equal a.e. in the sense of 
    \begin{equation*}
        \bm{X},\bm{Y} \colon \Omega \sto C([0,T],\R^n).
    \end{equation*}
    Therefore, $\bm{X} = \bm{Y}$ in $\mathcal{E}_T$ is equivalent to
    \begin{equation*}
        \Pb(\bm{X}_t(\omega) = \bm{Y}_t(\omega),~t \in [0,T]) = 1,
    \end{equation*}
    in other words, process $\bm{X}$ and $\bm{Y}$ are indistinguishable.
    \begin{itemize}
        \item \textbf{Check:} If $\bm{X}$ is a solution, then for any $T > 0$, $\bm{X} \in \mathcal{E}_T$.

        Fix $T > 0$. Define a stopping time
        \begin{equation*}
            \tau_n = \inf\bb{t \geq 0 \colon \norm{\bm{X}_t} > n}.
        \end{equation*}
        and consider $\bm{X}_{t \wedge \tau_n}$ for $ t \leq T$, i.e.,
        \begin{equation*}
            \bm{X}_{t\wedge \tau_n} = \bm{X}_0 + \int_0^{t \wedge \tau_n}b(s,\bm{X}_s)ds + \int_0^{t \wedge \tau_n} \sigma(s,\bm{X}_s)d\bm{W}_s.
        \end{equation*}
        Because $\norm{\bm{a} + \bm{b} + \bm{c}}^2 \leq 3(\norm{\bm{a}}^2 + \norm{\bm{b}}^2 + \norm{\bm{c}}^2)$,
        \begin{equation*}
            \norm{\bm{X}_{t\wedge \tau_n}}^2 \leq 3\bc{\norm{\bm{X}_0}^2+\norm{\int_0^{t \wedge \tau_n}b(s,\bm{X}_s)ds}^2 + \norm{\int_0^{t \wedge \tau_n} \sigma(s,\bm{X}_s)d\bm{W}_s}^2}.
        \end{equation*}
        It follows that
        \begin{align*}
            \E\bj{\sup_{0\leq u \leq t \wedge \tau_n}\norm{\bm{X}_{u}}^2} &\leq 3\left(\E\bj{\norm{\bm{X}_0}^2}+\E\bj{\sup_{0\leq u \leq t \wedge \tau_n}\norm{\int_0^{u}b(s,\bm{X}_s)ds}^2} \right. \\
            &\quad \left.+ \E\bj{\sup_{0\leq u \leq t \wedge \tau_n}\norm{\int_0^{u} \sigma(s,\bm{X}_s)d\bm{W}_s}^2}\right).
        \end{align*}
        First, by condition (\ref{eq:solcond1}),
        \begin{align*}
            \E\bj{\sup_{0\leq u \leq t \wedge \tau_n}\norm{\int_0^{u}b(s,\bm{X}_s)ds}^2} &\leq K^2\E\bj{\bc{\int_0^{t \wedge \tau_n}\bc{1+\norm{\bm{X}_s}}ds}^2} \\
            &\leq K^2(t \wedge \tau_n) \E\bj{\int_0^{t}\bc{1+\norm{\bm{X}_{s \wedge \tau_n}}}^2ds} \\
            &\leq 2K^2 T \int_0^{t} \bc{1 + \E\bj{\norm{\bm{X}_{s \wedge \tau_n}}^2}}ds \\
            &\leq  2K^2 T \int_0^{t} \bc{1 + \E\bj{\sup_{0\leq u \leq s \wedge \tau_n}\norm{\bm{X}_{u}}^2}}ds,
        \end{align*}
        where the second inequality is by
        \begin{equation*}
            \int^{t \wedge \tau_n}_0 g(X_s) ds \leq \int_0^t g(X_{s \wedge \tau_n})ds
        \end{equation*}
        for any positive $g$. Because
        \begin{equation*}
            \int_0^{t \wedge \tau_n} \sigma(s,\bm{X}_s)d\bm{W}_s
        \end{equation*}
        is a martingale, by Doob's maximal inequality,
        \begin{align*}
            \E\bj{\sup_{0\leq u \leq t}\norm{\int_0^{u  \wedge \tau_n} \sigma(s,\bm{X}_s)d\bm{W}_s}^2} &\leq 4\E\bj{\norm{\int_0^{t  \wedge \tau_n} \sigma(s,\bm{X}_s)d\bm{W}_s}^2} \\
            &= 4 \E\bj{\int_0^{t  \wedge \tau_n} \norm{\sigma(s,\bm{X}_s)}_{\op{F}}^2 ds} \\
            &\leq 4K^2 \E\bj{\int_0^{t  \wedge \tau_n} (1 + \norm{\bm{X}_s})^2 ds}\\
            &\leq 8K^2 \int_0^{t} \bc{1 + \E\bj{\sup_{0\leq u \leq s \wedge \tau_n}\norm{\bm{X}_{u}}^2}}ds.
        \end{align*}
        Therefore, we have
        \begin{equation*}
            \E\bj{\sup_{0\leq u \leq t \wedge \tau_n}\norm{\bm{X}_{u}}^2} \leq 3\left(
                \E\bj{\norm{\bm{X}_0}^2}+ 2\left(K^2 T+4 K^2\right)\int_0^{t} \bc{1 + \E\bj{\sup_{0\leq u \leq s \wedge \tau_n}\norm{\bm{X}_{u}}^2}}ds
            \right).
        \end{equation*}
        Note that $\bm{X}_0 = \bm{Z} \in L^2$, so $\E\bj{\norm{\bm{X}_0}^2} < \infty$. By applying Gr\"onwall's lemma to
        \begin{equation*}
            t \mapsto \E\bj{\sup_{0\leq u \leq t \wedge \tau_n}\norm{\bm{X}_{u}}^2},
        \end{equation*}
        we have
        \begin{equation*}
            \E\bj{\sup_{0\leq u \leq T \wedge \tau_n}\norm{\bm{X}_{u}}^2} < C,
        \end{equation*}
        for some $C$. As $n \sto \infty$, by Fatou's lemma,
        \begin{equation*}
            \E\bj{\sup_{0\leq u \leq T}\norm{\bm{X}_{u}}^2} < C,
        \end{equation*}
        and so $\bm{X} \in \mathcal{E}_t$.

        \item Fix $T = \delta$, for a small $\delta$. Let $\bm{X} \in \mathcal{E}_T$. Define
        \begin{equation*}
            \Phi(\bm{X})_t = \bm{Z} + \int_0^tb(s,\bm{X}_s)ds + \int_0^t \sigma(s,\bm{X}_s)d\bm{W}_s.
        \end{equation*}
        Note that it is obviously $\Phi(\bm{X})_t$ is $\mathcal{F}^W_t \vee \sigma(\bm{Z})$-measurable. In particular, by above we have
        \begin{equation*}
            \bm{\Phi(X)} \leq 3\bc{\E[\bm{Z}^2] + 2\left(K^2 T^2+4 K^2T\right)(1 + \norm{X}_E^2)} < \infty,
        \end{equation*}
        i.e., $\Phi \colon \mathcal{E}_T \sto \mathcal{E}_T$. Moreover, by applying condition (\ref{eq:solcond2}) and similarly reasoning as above, we have
        \begin{equation*}
            \|\Phi(\bm{X})-\Phi(\bm{Y})\|^2 \leq 2\left(K^2 T^2+4 K^2 T\right)\|\bm{X}-\bm{Y}\|^2.
        \end{equation*}
        When $T = \delta$ small enough, $\Phi$ is a contraction. So by Banach fixed point theorem, there exists a unique fixed point
        \begin{equation*}
            \bm{X}_t = \bm{Z} + \int_0^tb(s,\bm{X}_s)ds + \int_0^t \sigma(s,\bm{X}_s)d\bm{W}_s,
        \end{equation*}
        which provides a desired solution locally.

       \item In order to get a solution on $[0,\infty)$, we can apply previous step to get a solution on $[T_n,T_{n+1}]$, for $T_{n+1} - T_n$ sufficiently small. Then as $T_n \sto \infty$, we get the solution. Moreover, the uniqueness is because of the uniqueness of each interval $[T_n,T_{n+1}]$. \qedhere
    \end{itemize}
\end{proof}
\begin{rmk}
    In the proof, the truncation by $\tau_n$ is important, because in order to apply Doob's maximal inequality, we need the $L^2$-integrability. 
\end{rmk}
\begin{rmk}
    In fact, if $\bm{Z} \in L^p$, then we can get
    \begin{equation*}
        \E\bj{\sup_{0 \leq s \leq T}\norm{\bm{X}_s}^p} <\infty,
    \end{equation*}
    for all $p \geq 1$, by using by replacing Doob's inequality with Burkholder–Davis–Gundy (BDG) inequality on $\bm{M}_{t \wedge \tau_n} = \int_0^{t \wedge \tau_n} \sigma(s,\bm{X}_s)d\bm{W}_s$, which gets that
    \begin{equation*}
        \E\bj{\sup_{0 \leq u \leq t \wedge \tau_n} \norm{\int_0^{u} \sigma(s,\bm{X}_s)d\bm{W}_s}^p} \leq C_p\E\bj{\bc{\int_0^{t  \wedge \tau_n} \norm{\sigma(s,\bm{X}_s)}_{\op{F}}^2 ds}^{p/2}}.
    \end{equation*}
\end{rmk}

\begin{thm}[Burkholder-Davis-Gundy]
    Let $M$ be a continuous local martingale with $M_0 = 0$ and $\tau$ is a stopping time. Then for any $1\leq p < \infty$, there exists $c_p$ and $C_p$ such that
    \begin{equation*}
        c_p\E\bj{\inn{M}_\tau^{p/2}} \leq \E\bj{(M^*_\tau)^p} \leq C_p\E\bj{\inn{M}_\tau^{p/2}},
    \end{equation*}
    where $M^*_t = \sup_{0\leq s\leq t}\abs{M}_s$.
\end{thm}

If we only have the local Lipschitz continuity, then the solution also uniquely exists up to a stopping time.
\begin{defn}
    Consider a SDE
    \begin{equation*}
        d\bm{X}_t = b(t,\bm{X}_t)dt+\sigma(t,\bm{X}_t)d\bm{W}_t,\quad 0\leq t \leq T,~\bm{X}_0 = \bm{Z}.
    \end{equation*}
    Let $\tau$ be a $\mathbb{F}^W$-stopping time. $X$ is a solution up to $\tau$ if
    \begin{itemize}
        \item $\bm{X}_t$ defined on $\bb{t \leq \tau}$ is continuous and $\mathcal{F}^W_t \vee \sigma(\bm{Z})$-measurable.

        \item $\Pb\bc{\int_0^T \abs{b_i(s,\bm{X}_s)} + \abs{\sigma_{ij}(s,\bm{X}_s)}^2 ds < \infty} = 1$.

        \item $\bm{X}_{t} = \bm{Z} + \int_0^{t}b(s,\bm{X}_s)ds+\int_0^{t}\sigma(s,\bm{X}_s)d\bm{W}_s$ on $\bb{t < \tau}$ a.e..
    \end{itemize}
\end{defn}
\begin{rmk}
    Simply, we denote $\bm{X} = (\bm{X}_t)_{t < \tau}$. Also, we have another expression
    \begin{equation*}
        \bm{X}_{t \wedge \tau} = \bm{Z} + \int_0^{t \wedge \tau}b(s,\bm{X}_s)ds+\int_0^{t \wedge \tau}\sigma(s,\bm{X}_s)d\bm{W}_s.
    \end{equation*}
    Note that these two are actually equivalent. It seems the second one is stronger, because first for $\omega \in \bb{t < \omega}$, it has the above expression
    \begin{equation*}
        \bm{X}_{t}(\omega) = \bm{Z}(\omega) + \int_0^{t}b(s,\bm{X}_s(\omega))ds+\int_0^{t}\sigma(s,\bm{X}_s(\omega))d\bm{W}_s
    \end{equation*}
    but for $\omega \in \bb{t \geq \tau}$, it further requires
    \begin{equation*}
        \bm{X}_{\tau(\omega)}(\omega) = \bm{Z}(\omega) + \int_0^{\tau(\omega)}b(s,\bm{X}_s(\omega))ds+\int_0^{\tau(\omega)}\sigma(s,\bm{X}_s(\omega))d\bm{W}_s.
    \end{equation*}
    In fact, the first expression also has this formula by continuity.
\end{rmk}


\begin{thm}
    Fix $T > 0$ ($T$ could be $\infty$). Let $b \colon [0,T] \times \R^n \sto \R^n$ and $\sigma \colon [0,T] \times \R^n \sto \R^{n \times m}$ be locally Lipschitz, i.e., for every $N \geq 0$, there exists $C_N > 0$ such that
    \begin{equation*}
        \norm{b(t,x)-b(t,y)} + \norm{\sigma(t,x)-\sigma(t,y)}_{\op{F}} \leq C_N\norm{x-y},\quad \forall~B(\bm{0},N),~t \in [0,T].
    \end{equation*}
    Then for any $\bm{Z} \in L^2$, there exists a stopping time $\tau > 0$ a.e., and a solution $\bm{X}=(\bm{X}_t)_{t <\tau}$ up to $\tau$ for SDE
    \begin{equation*}
        \bm{X}_{t} = \bm{Z} + \int_0^{t}b(s,\bm{X}_s)ds+\int_0^{t}\sigma(s,\bm{X}_s)d\bm{W}_s,\quad t < \tau
    \end{equation*}
    This solution unique in the sense that if there exist another stopping time $\tau^\prime$ and a continuous process $\bm{X}^\prime$ such that
    \begin{equation*}
        \bm{X}^\prime_{t} = \bm{Z} + \int_0^{t}b(s,\bm{X}^\prime_s)ds+\int_0^{t}\sigma(s,\bm{X}^\prime_s)d\bm{W}_s,,\quad t < \tau^\prime,
    \end{equation*}
    then $\tau^\prime \leq \tau$ and 
    \begin{equation*}
        \bm{X}^\prime_t\mathbb{I}_{\bb{t < \tau^\prime}} = \bm{X}_t\mathbb{I}_{\bb{t < \tau^\prime}}.
    \end{equation*}
    Then $\tau = \tau(\bm{Z})$ is called the explosion time.
\end{thm}
\begin{proof}
    \begin{itemize}
        \item Fix $N > 0$. Define
        \begin{equation*}
            \pi_N(\bm{x}):= \begin{cases}\bm{x}, & \|\bm{x}\| \leq N \\ \frac{N}{\|\bm{x}\|} \bm{x}, & \|\bm{x}\|>N\end{cases}
        \end{equation*}
        and let
        \begin{equation*}
            b_N(t, \bm{x}):=b\left(t, \pi_N(\bm{x})\right), \quad \sigma_N(t, \bm{x}):=\sigma\left(t, \pi_N(\bm{x})\right) .
        \end{equation*}
        Then we have
        \begin{equation*}
            \left\|b_N(t, \bm{x})-b_N(t, \bm{y})\right\|=\left\|b\left(t, \pi_N(\bm{x})\right)-b\left(t, \pi_N(\bm{y})\right)\right\| \leq C_N\left\|\pi_N(\bm{x})-\pi_N(\bm{y})\right\| \leq C_N\|\bm{x}-\bm{y}\|
        \end{equation*}
        for any $t \in [0,T]$ and $\bm{x},\bm{y} \in \R^n$. Similarly, $\left\|\sigma_N(t, \bm{x})-\sigma_N(t, \bm{y})\right\| \leq C_N\|\bm{x}-\bm{y}\|$. Then by above theorem, there exists a unique solution $\bm{X}^N$ such that
        \begin{equation*}
            \bm{X}_t^N=\bm{Z}+\int_0^t b_N\left(s, \bm{X}_s^N\right) d s+\int_0^t \sigma_N\left(s, \bm{X}_s^N\right) d W_s, \quad t \in[0, T].
        \end{equation*}
        Now define
        \begin{equation*}
            \tau_N:=\inf \left\{t \geq 0:\left\|\bm{X}_t^N\right\| \geq N\right\} \wedge T .
        \end{equation*}
        Then on $\bb{t < \tau_N}$, because $\left\|\bm{X}_t^N\right\| < N$,
        \begin{equation*}
            b_N\left(t, \bm{X}_t^N\right)=b\left(t, \bm{X}_t^N\right), \quad \sigma_N\left(t, \bm{X}_t^N\right)=\sigma\left(t, \bm{X}_t^N\right) .
        \end{equation*}
        which implies that $\bm{X}^N$ is a solution of original SDE on $\bb{t < \tau_N}$.

        \item For $M > N$, we can have $\tau_M$ and $\bm{X}^M$. First, for all $\norm{\bm{x}} \leq N < M$, $\pi_N(\bm{x}) = \pi_M(\bm{x})$, hence
        \begin{equation*}
            b_N(t, x)=b_M(t, x)=b(t, x), \quad \sigma_N(t, x)=\sigma_M(t, x)=\sigma(t, x) .
        \end{equation*}
        Let
        \begin{equation*}
            \tau_N^M:=\inf \left\{t \geq 0:\left\|X_t^M\right\| \geq N\right\} \wedge T
        \end{equation*}
        On $t < \tau_M^N \wedge \tau_N$, because $\bm{X}^M,\bm{X}^N \in B(\bm{0},N)$, they solve the same SDE, So
        \begin{equation*}
            X_t^N=X_t^M, \quad \forall~ 0 \leq t \leq \tau_N \wedge \tau_N^M,
        \end{equation*}
        which implies that $\tau_N = \tau_N^M$. Moreover, it obvious that $\tau_N^M \leq \tau_M$. So $\bb{\tau_N}_N$ is an increasing sequence of stopping times. Let
        \begin{equation*}
            \tau = \lim_{N \sto \infty}\tau_N \in [0,\infty].
        \end{equation*}
        $\tau > 0$ because we can choose $N > \norm{\bm{Z}} + 1$.

        \item On $\omega \in \bb{t < \tau}$, choose $N$ such that $t < \tau_N(\omega)$, define
        \begin{equation*}
            \bm{X}_t(\omega) \defeq \bm{X}^N_t(\omega),
        \end{equation*} 
        which is well-defined because for $t < \tau_N \leq \tau_M$, $\bm{X}^N_t = \bm{X}^M_t$. Clearly, $\bm{X}_0 = \bm{Z}$ and on $\bb{t < \tau_N}$,
        \begin{equation*}
            \bm{X}_{t} = \bm{Z} + \int_0^{t}b(s,\bm{X}_s)ds+\int_0^{t}\sigma(s,\bm{X}_s)d\bm{W}_s,
        \end{equation*}
        i.e.,
        \begin{equation*}
            \bm{X}_{t \wedge \tau_N} = \bm{Z} + \int_0^{t \wedge \tau_N}b(s,\bm{X}_s)ds+\int_0^{t \wedge \tau_N}\sigma(s,\bm{X}_s)d\bm{W}_s.
        \end{equation*}
        As $N \sto \infty$,
        \begin{equation*}
            \bm{X}_{t \wedge \tau} = \bm{Z} + \int_0^{t \wedge \tau}b(s,\bm{X}_s)ds+\int_0^{t \wedge \tau}\sigma(s,\bm{X}_s)d\bm{W}_s.
        \end{equation*}

        \item Uniqueness: If we have $\bm{X}^\prime$ and $\tau^\prime$ such that
        \begin{equation*}
            \bm{X}_{t \wedge \tau^{\prime}}^{\prime}=Z+\int_0^{t \wedge \tau^{\prime}} b\left(s, \bm{X}_s^{\prime}\right) d s+\int_0^{t \wedge \tau^{\prime}} \sigma\left(s, \bm{X}_s^{\prime}\right) d \bm{W}_s
        \end{equation*}
        Let
        \begin{equation*}
            \rho_N:=\inf \left\{t \geq 0:\left\|X_t\right\| \geq N\right\}, \quad \rho_N^{\prime}:=\inf \left\{t \geq 0:\left\|X_t^{\prime}\right\| \geq N\right\}.
        \end{equation*}
        Then on $\bb{t < \rho_N \wedge \rho_N^{\prime} \wedge \tau \wedge \tau^{\prime}}$, both satisfies the original SDE. Define
        \begin{equation*}
            \bm{Y}_t:=\bm{X}_{t \wedge \rho_N \wedge \rho_N^{\prime} \wedge \tau \wedge \tau^{\prime}}-\bm{X}_{t \wedge \rho_N \wedge \rho_N^{\prime} \wedge \tau \wedge \tau^{\prime}}^{\prime},
        \end{equation*}
        and so 
        \begin{equation*}
            \bm{Y}_t=\int_0^{t \wedge \rho_N \wedge \rho_N^{\prime} \wedge \tau \wedge \tau^{\prime}}\left[b\left(s, \bm{X}_s\right)-b\left(s, \bm{X}_s^{\prime}\right)\right] d s+\int_0^{t \wedge \rho_N \wedge \rho_N^{\prime} \wedge \tau \wedge \tau^{\prime}}\left[\sigma\left(s, \bm{X}_s\right)-\sigma\left(s, \bm{X}_s^{\prime}\right)\right] d \bm{W}_s
        \end{equation*}
        By similar reasoning, for $f(t) \defeq \mathbb{E}\left(\sup _{0 \leq u \leq t}\left\|\bm{Y}_u\right\|^2\right)$,
        \begin{equation*}
            f(t) \leq K \int_0^tf(s)ds.
        \end{equation*}
        Because $f(t) \geq 0$ and $f(0) = 0$, $f(t) \equiv 0$, i.e.,
        \begin{equation*}
            \mathbb{E}\left(\sup _{0 \leq u \leq t}\left\|\bm{Y}_u\right\|^2\right) = 0~\quad~ \bm{Y}_t = 0,
        \end{equation*}
        which implies that
        \begin{equation*}
            \bm{X}_{t \wedge \rho_N \wedge \rho_N^{\prime} \wedge \tau \wedge \tau^{\prime}}=\bm{X}_{t \wedge \rho_N \wedge \rho_N^{\prime} \wedge \tau \wedge \tau^{\prime}}^{\prime}
        \end{equation*}
        Then as $N \sto \infty$,
        \begin{equation*}
            \bm{X}_t=\bm{X}_t^{\prime} \quad \forall~ t<\tau \wedge \tau^{\prime}.
        \end{equation*}
        Finally, we can show $\tau^{\prime} \leq \tau$. By continuity on $[0,\tau^\prime(\omega)]$ a.e.,
        \begin{equation*}
            \sup _{0 \leq s \leq \tau^{\prime}(\omega)}\left\|\bm{X}_s^{\prime}(\omega)\right\|<N .
        \end{equation*}
        Then by the uniqueness of solution,
        \begin{equation*}
            \bm{X}_t^{\prime}(\omega)=\bm{X}_t^N(\omega) \quad \forall~ t \leq \tau^{\prime}(\omega) .
        \end{equation*}
        Therefore,
        \begin{equation*}
            \sup _{0 \leq s \leq \tau^{\prime}(\omega)}\left\|\bm{X}_s^N(\omega)\right\|<N,
        \end{equation*}
        which implies that $ \tau(\omega) \geq \tau_N(\omega) > \tau^\prime(\omega)$ a.e.. \qedhere
    \end{itemize}
\end{proof}
\begin{rmk}
    Note that on our construction, we only care about the definition of $\bm{X}$ on $\bb{t < \tau}$. For $t \geq \tau$, we can extend it by adding a cemetery state $\Delta = \infty$,
    \begin{equation*}
        \tilde{\bm{X}}_t = \begin{cases}
            \bm{X}_t,& t < \tau \\
            \Delta,& t \geq \tau.
        \end{cases}
    \end{equation*}
    Therefore, we consider solution valued on $\R^n \cup {\infty}$, moreover, it is continuous by definition of $\tau$.
\end{rmk}

Moreover, a solution of SDE is a Markov process

\begin{thm}[Markov Property]
    Assume $b,\sigma$ are locally Lipschitz continuous. And $\bm{X}$ is a solution up to $\tau$. Let $f \colon \R^n \sto \R$ be a bounded Borel function. Then for any $t,s \geq 0$, we have
    \begin{equation*}
        \E_x\bj{f(\bm{X}_{t+s}) \mid \mathcal{F}^B_t} = \E_{\bm{X}_t}[f(\bm{X}_s)],
    \end{equation*}
    on $\bb{t < \tau}$.
\end{thm}

\begin{thm}[Strong Markov Property]
    Assume $b,\sigma$ are locally Lipschitz continuous. And $\bm{X}$ is a solution up to $\tau$. Let $f \colon \R^n \sto \R$ be a bounded Borel function and $T < \infty$ be a stopping time w.s.t. $\mathbb{F}^B$. Then
    \begin{equation*}
        \E_x\bj{f(\bm{X}_{T+s}) \mid \mathcal{F}^B_T} = \E_{\bm{X}_T}[f(\bm{X}_s)],
    \end{equation*}
    on $\bb{T < \tau}$.
\end{thm}

\section{Weak Solution}

\begin{defn}
    Given SDE,
    \begin{equation}\label{eq:sdeweakstrong}
        d\bm{X}_t = b(t,\bm{X}_t)dt+\sigma(t,\bm{X}_t)d\bm{W}_t.
    \end{equation}
    \begin{enumerate}[label=(\arabic{*})]
        \item A strong solution of (\ref{eq:sdeweakstrong}) on a give probability space $(\Omega,\mathcal{F},\Pb)$ and w.s.t. the fixed Brownian motion $\bm{W}$ and initial value $\bm{Z}$, is a stochastic process $\bm{X}$ with continuous paths and with the following properties:
        \begin{enumerate}[label=(\roman*)]
            \item $\bm{X}$ is adapted to $\mathbb{F}^W = (\mathcal{F}^W_t)_{t \geq 0}$,
            \item $\Pb(\bm{X}_0 = \bm{Z}) = 1$,
            \item $\Pb\bc{\int_0^T \abs{b_i(s,\bm{X}_s)} + \abs{\sigma_{ij}(s,\bm{X}_s)}^2 ds < \infty} = 1$, for all $i,j$,
            \item $\bm{X}$ satisfies the integral version
            \begin{equation*}
                \bm{X}_t = \bm{X}_0 + \int_0^tb(s,\bm{X}_s)ds+\int_0^t\sigma(s,\bm{X}_s)d\bm{W}_s.
            \end{equation*}
        \end{enumerate}

        \item A weak solution of (\ref{eq:sdeweakstrong}) is a triple $\bc{(\bm{X},\bm{W}),(\Omega,\mathcal{F},\mathbb{P}), \mathbb{F}=(\mathcal{F}_t)}$ such that
        \begin{enumerate}[label=(\roman*)]
            \item $(\Omega,\mathcal{F},\mathbb{P})$ is a probability space,
            \item $\mathbb{F}$ is a filtration on $(\Omega,\mathcal{F},\mathbb{P})$ with the usual condition,
            \item $\bm{X}$ is a continuous $\mathbb{F}$-adapted process,
            \item $\bm{W} = (\bm{W}_t,\mathcal{F}_t)$ is a Brownian motion,
            \item $\Pb\bc{\int_0^T \abs{b_i(s,\bm{X}_s)} + \abs{\sigma_{ij}(s,\bm{X}_s)}^2 ds < \infty} = 1$, for all $i,j$,
            \item $\bm{X}$ satisfies the integral version
            \begin{equation*}
                \bm{X}_t = \bm{X}_0 + \int_0^tb(s,\bm{X}_s)ds+\int_0^t\sigma(s,\bm{X}_s)d\bm{W}_s.
            \end{equation*}
        \end{enumerate}
    \end{enumerate}
\end{defn}
\begin{rmk}
    It is obviously that the existence of strong solution implies that the solution is also a weak solution. But the existence of weak solution does not implies the existence of strong solution.
\end{rmk}

\begin{exam}[Tanaka Equation]
    Consider the SDE
    \begin{equation*}\label{eq:tanaka}
        dX_t = \sgn(X_t)dW_t.
    \end{equation*}
    Note that $\sigma(t,x) = \sgn(x)$ does not satisfy the Lipschitz condition.

    \noindent \textbf{Claim:} (\ref{eq:tanaka}) has no strong solution.

    \noindent Suppose (\ref{eq:tanaka}) has a strong solution $X$. Then
    \begin{equation*}
        X_t = \int_0^t \sgn(X_s)dW_s~\Rightarrow~ \inn{X}_t = t.
    \end{equation*}
    By L\'evy theorem, $X$ is a Brownian motion. On the other hand,
    \begin{equation*}
        dW_t = \sgn(X_t)dX_t~\Rightarrow~ W_t = \int_0^t \sgn{X_s}ds,
    \end{equation*}
    which means $W$ is a Brownian motion w.s.t. $\mathbb{F}^X$, i.e., $\mathbb{F}^W \subset \mathbb{F}^X$. By the Tanaka equation,
    \begin{equation*}
        W_t = \abs{X_t} - L^X_t,
    \end{equation*}
    which implies that $\mathcal{F}^W_t \neq \mathcal{F}^X_t$, contradicting to $X$ adapted to $\mathbb{F}^W$.

    \noindent (\ref{eq:tanaka}) has a weak solution. Choose $B = (B_t)_{t \geq 0}$ be a Brownian motion. Define
    \begin{equation*}
        \tilde{W}_t = \int_0^t\sgn(B_u)dB_u,
    \end{equation*}
    which is a Brownian motion. Then let $X=B$,
    \begin{equation*}
        d\tilde{W}_t = \sgn(X_t)dX_t~\Rightarrow~dX_t = \sgn(X_t)d\tilde{W}_t.
    \end{equation*}
\end{exam}

\begin{defn}
    For a given SDE,
    \begin{equation*}
        d\bm{X}_t = b(t,\bm{X}_t)dt+\sigma(t,\bm{X}_t)d\bm{W}_t, \quad \bm{X}_0 \sim \mu,
    \end{equation*}
    \begin{itemize}
        \item Weak uniqueness means that any two weak solutions has the same law.
        \item Pathwise uniqueness means that when fix $(\Omega,\mathcal{F},\mathbb{P}), \mathbb{F}=(\mathcal{F}_t)$, Brownian motion $\bm{W} = (\bm{W}_t,\mathcal{F}_t)$, and $\bm{X}_0$, two solutions are indistinguishable, i.e. strong uniqueness.
    \end{itemize}
\end{defn}

\begin{thm}[Yamada–Watanabe]
    If both weak existence and pathwise uniqueness hold, then weak uniqueness also holds.
\end{thm}
\begin{cor}
    If $b,\sigma$ are locally Lipschitz continuous, then the solution is weakly unique.
\end{cor}



\section{Feynman-Kac Formula}

\begin{thm}[Feynman-Kac Formula]
    Consider SDE
    \begin{equation*}
        dX_t = \beta(t,X_t)dt+ \sigma(t,X_t)dW_t.
    \end{equation*}
    Let $f$ be a Borel-measurable function. Fix $T > 0$ and $t \in [0,T]$. Define
    \begin{equation*}
        g(t,x) = \E\bj{f(X_T) \mid X_t = x} = \E^{t,x}[f(X_T)].
    \end{equation*}
    Assume $g(t,x) < \infty$. Then $g(t,x)$ satisfies PDE
    \begin{equation*}
        \frac{\partial}{\partial t}g + \beta \frac{\partial}{\partial x} g + \frac{1}{2} \sigma^2 \frac{\partial^2}{\partial x^2}g = 0
    \end{equation*}
    with terminal $g(T,x) = f(x)$.
\end{thm}
\begin{rmk}
    Note that $(g(t,X_t))_{0 \leq t \leq T}$ is a martingale, because by It\^o formula
    \begin{align*}
        dg(t,X_t) &= g_t(t,X_t)dt + g_x(t,X_t)dX_t + \frac{1}{2}g_{xx}(t,X_t)(dX_t)^2 \\
        &= g_tdt+g_x(\beta dt + \sigma dW_t) + \frac{1}{2}g_{xx} \sigma^2 dt \\
        &= g_x \sigma dW_t + \bc{g_t + \beta g_x + \frac{1}{2}\sigma^2 g_{xx}}dt \\
        &= g_x \sigma dW_t.
    \end{align*}
\end{rmk}

\begin{thm}[Discounted Feynman-Kac Formula]
    Consider SDE
    \begin{equation*}
        dX_t = \beta(t,X_t)dt+ \sigma(t,X_t)dW_t.
    \end{equation*}
    Let $f$ be a Borel-measurable function and $r$ be a constant. Fix $T > 0$ and $t \in [0,T]$. 
    Define
    \begin{equation*}
        h(t,x) = \E^{t,x}\bj{e^{-r(T-t)}f(X_T)}.
    \end{equation*}
    Then $h$ satisfies
    \begin{equation*}
        h_t(t,x) + \beta(t,x)h_x(t,x)+ \frac{1}{2}\sigma(t,x)h_{xx}(t,x) = r h(t,x)
    \end{equation*}
    with terminal $h(T,x) = f(x)$.
\end{thm}

\section{Discretization}

\paragraph{Euler–Maruyama Algorithm.} For a given SDE,
\begin{equation*}
    d\bm{X}_t = b(t,\bm{X}_t)dt+\sigma(t,\bm{X}_t)d\bm{W}_t, \quad t \in [0,T],\quad \bm{X}_0 = \bm{Z},
\end{equation*}
with globally Lipschitz continuous $b,\sigma$. Let $h = T/N$ for a large $N \in \N$ and $t_k = kh$ for $k = 0,1,\cdots,N$. Construction $\{\bm{X}^h_k\}_{k=0}^N$ by setting $\bm{X}^h_0 = \bm{Z} \in L^2$ and
\begin{equation*}
    \bm{X}^h_{k+1} = \bm{X}^h_k + b(t_k,\bm{X}^h_k)h + \sigma(t_k,\bm{X}^h_k)\Delta \bm{W}_k.
\end{equation*}
Note that $\Delta \bm{W}_k$ is replaced by $\sqrt{h}\bm{\xi}$ for $\bm{\xi} \sim \mathcal{N}(0,I)$ in practice.
\begin{thm}[Strong Convergence of EM]
    Using notations as above, there exists a constant $C>0$ such that
    \begin{equation*}
        \bc{\E\bj{\max_{0\leq k \leq N}\norm{\bm{X}_{t_k} - \bm{X}^h_k}^2}}^{1/2} \leq Ch^{1/2}.
    \end{equation*}
\end{thm}
\begin{proof}
    First, from the proof of the existence and uniqueness of solution, we have
    \begin{equation*}
        \sup_{0 \leq t \leq T}\E\bj{\norm{\bm{X}_t}^2} < \infty,\quad \max_{0\leq k \leq N} \E\bj{\norm{\bm{X}^h_{k}}^2} < \infty.
    \end{equation*}
    More, by the similar reasoning, there exists $C_1$ such for $s \in [t_k,t_{k+1}]$,
    \begin{equation*}
        \E\bj{\norm{\bm{X}_s - \bm{X}_{t_k}}^2} \leq C_1h,
    \end{equation*}
    when $h$ is small.

    Let $\bm{e}_k \defeq \bm{X}_{t_k} - \bm{X}^h_{k}$. Note that
    \begin{equation*}
        \bm{X}_{t_{k+1}}=\bm{X}_{t_k}+\int_{t_k}^{t_{k+1}} b\left(s,\bm{X}_s\right) d s+\int_{t_k}^{t_{k+1}} \sigma\left(s,\bm{X}_s\right) d \bm{W}_s
    \end{equation*}
    and $\bm{X}^h_{k+1} = \bm{X}^h_k + b(t_k,\bm{X}^h_k)h + \sigma(t_k,\bm{X}^h_k)\Delta \bm{W}_k$. Therefore, $\bm{e}_{k+1} = \bm{e}_k + \bm{A}_k + \bm{M}_k$, where
    \begin{align*}
        \bm{A}_k & :=\int_{t_k}^{t_{k+1}}\left(b\left(s,\bm{X}_s\right)-b\left(t_k,\bm{X}_k^h\right)\right) d s \\
        \bm{M}_k & :=\int_{t_k}^{t_{k+1}}\left(\sigma\left(s,\bm{X}_s\right)-\sigma\left(t_k,\bm{X}_k^h\right)\right) d \bm{W}_s.
    \end{align*}
    Therefore, by $K$-Lipschitz continuity,
    \begin{align*}
        \E\bj{\norm{\bm{A}_k}^2} &= \E\bj{\norm{\int_{t_k}^{t_{k+1}}\left(b\left(s,\bm{X}_s\right)-b\left(t_k,\bm{X}_k^h\right)\right) d s }^2} \\
        &\leq K^2\E\bj{\bc{\int_{t_k}^{t_{k+1}} \norm{\bm{X}_s - \bm{X}_h^k} ds}^2} \\
        &\leq K^2h\int_{t_k}^{t_{k+1}} \E\bj{\norm{\bm{X}_s - \bm{X}_h^k}^2} ds \\
        &\leq K^2h\int_{t_k}^{t_{k+1}} \E\bj{\norm{\bm{X}_s - \bm{X}_{t_k}}^2} + \E\bj{\norm{\bm{e}_k}^2} ds \\
        &\leq K^2h^2(C_1h + \E\bj{\norm{\bm{e}_k}^2}),
    \end{align*}
    and similarly,
    \begin{equation*}
        \E\bj{\norm{\bm{M}_2}^2} = \int_{t_k}^{t_{k+1}} \E\norm{\sigma\left(s,\bm{X}_s\right)-\sigma\left(t_k,\bm{X}_k^h\right)}^2 ds \leq K^2h(C_1h + \E\bj{\norm{\bm{e}_k}^2}),
    \end{equation*}
    Therefore,
    \begin{align*}
        \E\bj{\norm{\bm{e}_{k+1}}^2} &\leq 3\bc{\E\bj{\norm{\bm{e}_k}^2} + \E\bj{\norm{\bm{A}_k}^2} +\E\bj{\norm{\bm{M}_k}^2}} \\
        &\leq 3(1 + C_2h)\E\bj{\norm{\bm{e}_k}^2} + C_3h^2,
    \end{align*}
    with $\E\bj{\norm{\bm{e}_{0}}^2} = 0$. By discrete Gr\"onwall,
    \begin{equation*}
        \E\bj{\norm{\bm{e}_k}^2} \leq C_3h^2\sum_{j=0}^{k-1}(1+C_2h)^{k-1-j} \leq C_3h^2\frac{(1+C_2h)^k-1}{(1+C_2h)-1} \leq Ch. \qedhere
    \end{equation*}
\end{proof}

\paragraph{Weak Convergence.} For a given SDE,
\begin{equation*}
    d\bm{X}_t = b(t,\bm{X}_t)dt+\sigma(t,\bm{X}_t)d\bm{W}_t,
\end{equation*}
we construct a discrete $(\bm{X}^h_k)_{k=0}^N$ with $h = T /N < 1$ and $t_k = kh$, which is said a weak approximation if 
\begin{equation*}
    \E[\varphi(\bm{X}^h_N)] \sto \E[\varphi(\bm{X}_T)].
\end{equation*}

Consider the following example. Let $b(t,x) = \mu(x)$ and $\sigma(t,x) = \sigma(x)$. Assume there exists a family of probability measures $\bb{\nu_x}_{x\in \R^d}$ such that
\begin{equation*}
    \mu(x) = \E_{\bm{Y} \sim \nu_x}[\bm{Y}],\quad \sigma(x)\sigma(x)^\top = \op{Cov}_{\bm{Y} \sim \nu_x}[\bm{Y}].
\end{equation*}
Let $(P_t)_{t \in [0,T]}$ be the corresponding semigroup. Note that the corresponding generator is given by
\begin{equation*}
    Af(x) = \sum_i \mu_i(x) \frac{\partial f}{\partial x^i}(x) + \frac{1}{2}\sum_{i,j}a_{ij}(x)\frac{\partial^2 f}{\partial x^i\partial x^j}(x) = \inn{\mu,\nabla f} + \frac{1}{2}a : D^2\varphi,
\end{equation*}
where $a(x) = \sigma(x)\sigma(x)^\top$; see in the next chapter.

We construct $\bm{X}^h_k$ by the following rule: first, drawing $\bm{V}_k$ such that
\begin{equation*}
    \bm{V}_k \mid \bm{X}^h_k = x ~\sim~\nu_x,
\end{equation*}
and then updating
\begin{equation*}
    \bm{X}^h_{k+1} = \bm{X}^h_k + \mu(\bm{X}^h_k)h + \sqrt{h}(\bm{V}_k - \mu(\bm{X}^h_k)).
\end{equation*}
Let $(R^h_k)_{k=0}^N$ be the corresponding semigroup.

\begin{thm}\label{thm:weakconv}
    Using notations as above, assume $\mu_i,\sigma_{ij} \in C^3_b(\R^n)$ and $a(x)$ is uniformly elliptic and regular enough. Further assume that $\sup_x\E_{\bm{Y} \sim \nu_x}\bj{\norm{\bm{Y}}^3}  < \infty$. Then there exists a $C_T > 0$ such that
    \begin{equation*}
        \norm{R^h_Nf - P_Tf}_\infty \leq C_Th^{1/2},\quad \forall~f \in C^3_b(\R^n)
    \end{equation*}
\end{thm}
\begin{lem}\label{lem:onesteptrans}
    Let $S_h = R^h_1$. For any $\varphi \in C^3_b$, 
    \begin{equation*}
        \sup_x\left|\left(S_h \varphi\right)(x)-\varphi(x)-h(A \varphi)(x)\right| \leq C h^{3 / 2}.
    \end{equation*}
\end{lem}
\begin{proof}
    Fix $x \in \R^n$. Let $\eta_x = V - \mu(x)$ for $V \sim \nu_x$. Then by definition,
    \begin{equation*}
        \mathbb{E}\left[\eta_x\right]=0, \quad \mathbb{E}\left[\eta_x \eta_x^{\top}\right]=a(x), \quad \mathbb{E}\left[\left|\eta_x\right|^3\right] \leq C .
    \end{equation*}
    Let
    \begin{equation*}
        \Delta:=\Delta(x, h):=\mu(x) h+\sqrt{h} \eta_x .
    \end{equation*}
    We have
    \begin{equation*}
        \left(S_h \varphi\right)(x)=\mathbb{E}[\varphi(x+\Delta)] .
    \end{equation*}
    (because $\E[\varphi(X,Y) \mid X = x] = \E[\varphi(x,Z)]$ for $Z \sim Y \mid X=x$). By Taylor expansion,
    \begin{equation*}
        \varphi(x+\Delta)=\varphi(x)+\nabla \varphi(x) \cdot \Delta+\frac{1}{2} \Delta^{\top} D^2 \varphi(x) \Delta+R(x, \Delta),
    \end{equation*}
    where
    \begin{equation*}
        |R(x, \Delta)| \leq C_{\varphi}|\Delta|^3,
    \end{equation*}
    by $\sup_x \abs{D^3\varphi(x)} < C_{\varphi}$. Therefore,
    \begin{equation*}
        \left(S_h \varphi\right)(x)-\varphi(x)=\nabla \varphi(x) \cdot \mathbb{E}[\Delta]+\frac{1}{2} \mathbb{E}\left[\Delta^{\top} D^2 \varphi(x) \Delta\right]+\mathbb{E}[R(x, \Delta)].
    \end{equation*}
    For the three terms on the RHS,
    \begin{itemize}
        \item $\mathbb{E}[\Delta]=\mu(x) h+\sqrt{h} \mathbb{E}\left[\eta_x\right]=\mu(x) h $.
        \item because $\mathbb{E}\left[\Delta \Delta^{\top}\right]= h a(x)+h^2 \mu(x) \mu(x)^{\top}$,
        \begin{equation*}
            \frac{1}{2} \mathbb{E}\left[\Delta^{\top} D^2 \varphi(x) \Delta\right]=\frac{1}{2}\left(h a(x)+h^2 \mu \mu^{\top}\right): D^2 \varphi(x)=\frac{1}{2} h a(x): D^2 \varphi(x)+O\left(h^2\right)
        \end{equation*}
        because $D^2\varphi$ is uniformly bounded.

        \item by $|a+b|^3 \leq C_1\left(|a|^3+|b|^3\right)$,
        \begin{equation*}
            |\Delta|^3 \leq C\left(h^3|\mu(x)|^3+h^{3 / 2}\left|\eta_x\right|^3\right) .
        \end{equation*}
        Since $\mu$ and $\E[\left|\eta_x\right|^3]$ are uniformly bounded,
        \begin{equation*}
            \mathbb{E}|\Delta|^3 \leq C_2\left(h^3+h^{3 / 2}\right) \leq C_3 h^{3 / 2} .
        \end{equation*}
        for $h < 1$.
    \end{itemize}
    Then we have
    \begin{equation*}
        \left(S_h \varphi\right)(x)-\varphi(x)=h(A \varphi)(x)+O\left(h^2\right)+O\left(h^{3 / 2}\right). \qedhere
    \end{equation*}
\end{proof}
\begin{lem}\label{lem:discreteu}
    Fix $f \in C^3_b(\R^n)$. Let $u(t,x) = \E[f(X_{T-t}) \colon X_0 = x] = P_{T-t}f(x)$ and $u_k(x) = u(t_k,x)$. There exists a $C^\prime$ such that
    \begin{equation*}
        \sup_x\left|u_k(x)-u_{k+1}(x)-h A u_{k+1}(x)\right| \leq C^\prime h^2 .
    \end{equation*}
\end{lem}
\begin{proof}
    Note that by the backward Kolmogorov equation
    \begin{equation*}
        \partial_tu(t,x) + Au(t,x) = 0,\quad u(T,x) = f(x).
    \end{equation*}
    Therefore,
    \begin{equation*}
        u_k(x)-u_{k+1}(x)=\int_{t_k}^{t_{k+1}} A u(s, x) d s,
    \end{equation*}
    which implies that
    \begin{equation*}
        u_k(x)-u_{k+1}(x)-h A u_{k+1}(x)=\int_{t_k}^{t_{k+1}}\left(A u(s, x)-A u\left(t_{k+1}, x\right)\right) d s.
    \end{equation*}
    Moreover, our ``regular enough'' for $a$ means that $u \in C_b^{1,3}([0,T] \times \R^n)$, so
    \begin{equation*}
        \left|A u(s, x)-A u\left(t_{k+1}, x\right)\right| \leq C_4\left|s-t_{k+1}\right|
    \end{equation*}
    by the boundedness of $\mu,\sigma$. Therefore,
    \begin{equation*}
        \abs{u_k(x)-u_{k+1}(x)-h A u_{k+1}(x)} \leq \int_{t_k}^{t_{k+1}} C_4\left|s-t_{k+1}\right| d s \leq C^\prime h^2. \qedhere
    \end{equation*}
\end{proof}
\begin{proof}[Proof of Theorem \ref{thm:weakconv}]
    Fix $f \in C^3_b(\R^n)$. Let $u(t,x) = \E[f(X_{T-t}) \colon X_0 = x] = P_{T-t}f(x)$ and $u_k = u(t_k,\cdot)$. We need to measure $\norm{R^h_Nf - u(0,\cdot)}$.
    \begin{itemize}
        \item First, let
        \begin{equation*}
            \delta_k(x):=\left(S_h u_{k+1}\right)(x)-u_k(x) .
        \end{equation*}
        Then
        \begin{equation*}
            \delta_k(x)=\left(S_h u_{k+1}(x)-u_{k+1}(x)-h A u_{k+1}(x)\right)+\left(u_{k+1}(x)+h A u_{k+1}(x)-u_k(x)\right) .
        \end{equation*}
        For the first term, by Lemma \ref{lem:onesteptrans},
        \begin{equation*}
            \left|S_h u_{k+1}(x)-u_{k+1}(x)-h A u_{k+1}(x)\right| \leq C h^{3 / 2} .
        \end{equation*}
        By Lemma \ref{lem:discreteu},
        \begin{equation*}
            \left|u_{k+1}(x)+h A u_{k+1}(x)-u_k(x)\right| \leq C^\prime h^2 .
        \end{equation*}
        Therefore,
        \begin{equation*}
            \left|\delta_k(x)\right| \leq C_5 h^{3 / 2} \quad \Longrightarrow \quad\left\|\delta_k\right\|_{\infty} \leq C_5 h^{3 / 2} .
        \end{equation*}

        \item Note that $u(T,x) = f(x)$ and $R^h_Nf(x) = \E\bj{f(\bm{X}^h_N) \mid X_0 = x} = \E_x[u(T,\bm{X}^h_N)]$. So
        \begin{equation*}
            R_N^h f(x)-P_T f(x)=\mathbb{E}_x\left[u\left(T, \bm{X}_N^h\right)-u(0, x)\right]=\sum_{k=0}^{N-1} \mathbb{E}_x\left[u\left(t_{k+1}, \bm{X}_{k+1}^h\right)-u\left(t_k, \bm{X}_k^h\right)\right] .
        \end{equation*}
        For each $k$, we have
        \begin{equation*}
            \mathbb{E}_x\left[u\left(t_{k+1}, \bm{X}_{k+1}^h\right)-u\left(t_k, \bm{X}_k^h\right)\right]=\mathbb{E}_x\left[\mathbb{E}\left[u_{k+1}\left(\bm{X}_{k+1}^h\right)-u_k\left(\bm{X}_k^h\right) \mid \bm{X}_k^h\right]\right].
        \end{equation*}
        Because 
        \begin{equation*}
            \mathbb{E}\left[u_{k+1}\left(\bm{X}_{k+1}^h\right) \mid \bm{X}_k^h=y\right]=\left(S_h u_{k+1}\right)(y),
        \end{equation*}
        we get
        \begin{equation*}
            \mathbb{E}_x\left[u\left(t_{k+1}, \bm{X}_{k+1}^h\right)-u\left(t_k, \bm{X}_k^h\right)\right]=\mathbb{E}_x\left[\delta_k\left(\bm{X}_k^h\right)\right],
        \end{equation*}
        which implies that
        \begin{equation*}
            \left\|R_N^h f-P_T f\right\|_{\infty}=\sup _x\left|\sum_{k=0}^{N-1} \mathbb{E}_x\left[\delta_k\left(\bm{X}_k^h\right)\right]\right| \leq \sum_{k=0}^{N-1} \sup _x\left|\mathbb{E}_x\left[\delta_k\left(\bm{X}_k^h\right)\right]\right| .
        \end{equation*}
        Moreover, since $\left|\mathbb{E}_x\left[\delta_k\left(\bm{X}_k^h\right)\right]\right| \leq\left\|\delta_k\right\|_{\infty}$,
        \begin{equation*}
            \left\|R_N^h f-P_T f\right\|_{\infty} \leq \sum_{k=0}^{N-1}\left\|\delta_k\right\|_{\infty} \leq \sum_{k=0}^{N-1} C_5 h^{3 / 2}=N C_5 h^{3 / 2}=\frac{T}{h} C_5 h^{3 / 2}=C_T h^{1 / 2} . \qedhere
        \end{equation*}
    \end{itemize}
\end{proof}








